\chapter{Sprint 5 : Développement du module d'Intelligence Artificielle (Prédiction multi-modèles)}
\label{chap:sprint5-ia}

\section*{Introduction}
\addcontentsline{toc}{section}{Introduction}

L'intelligence artificielle constitue la brique la plus innovante de notre solution décisionnelle. Alors que les tableaux de bord permettent de visualiser le passé et le présent, l'IA ouvre la voie à l'anticipation : elle permet de prédire les performances futures des agences, d'identifier les tendances émergentes et d'alerter sur les risques potentiels. Ce cinquième sprint est dédié au développement d'un module de prédiction multi-modèles, qui comparera plusieurs algorithmes pour offrir aux responsables de la BTL des prévisions fiables et éclairées.

Dans ce chapitre, nous commencerons par définir les objectifs précis du sprint, puis nous présenterons la conception à travers des diagrammes UML (cas d'utilisation, classes, séquence). Nous détaillerons ensuite la définition du problème de prédiction, la préparation des données, l'implémentation des différents modèles, leur évaluation et leur comparaison, ainsi que la sauvegarde et l'intégration dans l'API FastAPI. Enfin, nous aborderons l'interface React de visualisation des prédictions, les tests effectués et les livrables de ce sprint.

\section{Objectifs du Sprint}
\label{sec:sprint5-objectifs}

Le sprint 5 a pour objectifs principaux :

\begin{itemize}
    \item Définir le \textbf{problème de prédiction} en identifiant les variables cibles (dépôts, crédits, PNB) et les horizons temporels.
    \item Préparer les \textbf{données} pour l'IA : feature engineering, nettoyage, normalisation et division en ensembles d'entraînement et de test.
    \item Implémenter \textbf{plusieurs modèles de prédiction} avec Python et Scikit-learn : Régression Linéaire, Arbre de Décision, Forêt Aléatoire, Gradient Boosting et Réseau de Neurones.
    \item \textbf{Évaluer et comparer} les performances des modèles à l'aide de métriques standards (RMSE, MAE, R²).
    \item \textbf{Sauvegarder} les modèles entraînés pour une réutilisation future.
    \item \textbf{Intégrer} les modèles dans l'API FastAPI via un endpoint de prédiction.
    \item Développer une \textbf{interface React} permettant de sélectionner un modèle et de visualiser les prédictions.
\end{itemize}

À la fin de ce sprint, les responsables de la BTL pourront anticiper les performances futures de leurs agences et comparer la fiabilité des différents modèles de prédiction.

\section{Conception du Sprint}
\label{sec:sprint5-conception}

La phase de conception est essentielle pour structurer le module d'IA et définir ses composants. Elle permet de visualiser les interactions entre les acteurs et le système, de définir la structure des classes Python et de modéliser les flux d'exécution.

\subsection{Diagramme de cas d'utilisation}
\label{subsec:sprint5-use-case}

Le diagramme de cas d'utilisation présente les fonctionnalités offertes par le module de prédiction du point de vue des acteurs. Les acteurs identifiés sont :

\begin{itemize}
    \item \textbf{Directeur} : Peut utiliser les modèles de prédiction pour toutes les agences, comparer les résultats et identifier les tendances globales.
    \item \textbf{Responsable d'agence} : Peut utiliser les modèles de prédiction pour sa propre agence uniquement, afin d'anticiper ses performances futures.
    \item \textbf{Administrateur} : Peut gérer les modèles (ajout, suppression, mise à jour) et superviser leur performance.
    \item \textbf{Système} : Le module d'IA interagit avec le Data Warehouse pour récupérer les données historiques et avec l'API FastAPI pour exposer les prédictions.
\end{itemize}

\begin{figure}[H]
\centering
\includegraphics[scale=0.5]{images/diagramme_use_case_ia.png}
\caption{Diagramme de cas d'utilisation - Module de prédiction IA}
\label{fig:use-case-ia}
\end{figure}

\textbf{Description :} Le diagramme de cas d'utilisation présente les interactions entre les acteurs (Directeur, Responsable d'agence, Administrateur, Système) et les fonctionnalités du module de prédiction : entraînement des modèles, obtention de prédictions, comparaison des modèles, visualisation des résultats, gestion des modèles et consultation des métriques.

Les cas d'utilisation identifiés pour ce sprint sont :

\begin{table}[H]
\centering
\begin{tabular}{|p{4cm}|p{3.5cm}|p{6cm}|}
\hline
\textbf{Cas d'utilisation} & \textbf{Acteurs} & \textbf{Description} \\
\hline
Entraîner les modèles & Administrateur & Lancer l'entraînement des modèles sur les données historiques. \\
\hline
Obtenir une prédiction & Directeur, Responsable & Recevoir une prédiction pour une agence donnée. \\
\hline
Comparer les modèles & Directeur & Visualiser les performances relatives des différents modèles. \\
\hline
Visualiser les résultats & Directeur, Responsable & Consulter les prédictions sous forme graphique et tabulaire. \\
\hline
Gérer les modèles & Administrateur & Ajouter, mettre à jour ou supprimer des modèles. \\
\hline
Consulter les métriques & Administrateur & Visualiser les métriques de performance des modèles. \\
\hline
\end{tabular}
\caption{Cas d'utilisation du sprint 5}
\label{tab:use-cases-ia}
\end{table}

\subsection{Diagramme de classes}
\label{subsec:sprint5-classes}

Le diagramme de classes modélise la structure statique du module d'IA, en représentant les classes Python principales, leurs attributs et leurs méthodes.

Les principales classes identifiées sont :

\begin{itemize}
    \item \textbf{DataPreprocessor} : Classe responsable de la préparation des données.
        \begin{itemize}
            \item Attributs : data (DataFrame), features (liste), target (nom de la variable cible).
            \item Méthodes : clean\_data(), normalize\_data(), encode\_categorical(), split\_train\_test().
        \end{itemize}
    
    \item \textbf{BaseModel} : Classe abstraite définissant l'interface commune à tous les modèles.
        \begin{itemize}
            \item Méthodes : train(X\_train, y\_train), predict(X\_test), evaluate(X\_test, y\_test), save(path), load(path).
        \end{itemize}
    
    \item \textbf{LinearRegressionModel} : Modèle de régression linéaire (baseline).
        \begin{itemize}
            \item Hérite de BaseModel.
            \item Utilise LinearRegression de Scikit-learn.
        \end{itemize}
    
    \item \textbf{DecisionTreeModel} : Modèle d'arbre de décision.
        \begin{itemize}
            \item Hérite de BaseModel.
            \item Utilise DecisionTreeRegressor de Scikit-learn.
            \item Paramètres : max\_depth, min\_samples\_split.
        \end{itemize}
    
    \item \textbf{RandomForestModel} : Modèle de forêt aléatoire.
        \begin{itemize}
            \item Hérite de BaseModel.
            \item Utilise RandomForestRegressor de Scikit-learn.
            \item Paramètres : n\_estimators, max\_depth, min\_samples\_split.
        \end{itemize}
    
    \item \textbf{GradientBoostingModel} : Modèle de gradient boosting.
        \begin{itemize}
            \item Hérite de BaseModel.
            \item Utilise XGBoost ou LightGBM.
            \item Paramètres : learning\_rate, n\_estimators, max\_depth.
        \end{itemize}
    
    \item \textbf{NeuralNetworkModel} : Modèle de réseau de neurones simple.
        \begin{itemize}
            \item Hérite de BaseModel.
            \item Utilise MLPRegressor de Scikit-learn.
            \item Paramètres : hidden\_layer\_sizes, activation, max\_iter.
        \end{itemize}
    
    \item \textbf{ModelEvaluator} : Classe responsable de l'évaluation et de la comparaison des modèles.
        \begin{itemize}
            \item Méthodes : compute\_metrics(y\_true, y\_pred), compare\_models(models, X\_test, y\_test), generate\_comparison\_table().
        \end{itemize}
    
    \item \textbf{ModelSaver} : Classe responsable de la sauvegarde et du chargement des modèles.
        \begin{itemize}
            \item Méthodes : save\_model(model, path), load\_model(path), list\_available\_models().
        \end{itemize}
    
    \item \textbf{PredictionService} : Service intégré dans FastAPI pour exposer les prédictions.
        \begin{itemize}
            \item Méthodes : predict(agence\_id, model\_name, periode), get\_available\_models(), get\_model\_metrics(model\_name).
        \end{itemize}
\end{itemize}

\begin{figure}[H]
\centering
\includegraphics[scale=0.5]{images/diagramme_classes_ia.png}
\caption{Diagramme de classes - Module d'IA}
\label{fig:class-diagram-ia}
\end{figure}

\textbf{Description :} Le diagramme de classes présente la structure du module d'IA avec la classe abstraite BaseModel dont héritent les cinq modèles (LinearRegressionModel, DecisionTreeModel, RandomForestModel, GradientBoostingModel, NeuralNetworkModel). Les classes DataPreprocessor, ModelEvaluator, ModelSaver et PredictionService complètent l'architecture.

Les relations entre les classes sont les suivantes :

\begin{itemize}
    \item \textbf{DataPreprocessor} est utilisé par tous les modèles pour préparer les données.
    \item Chaque modèle (\textbf{LinearRegressionModel}, etc.) hérite de \textbf{BaseModel}.
    \item \textbf{ModelEvaluator} prend une liste de modèles en entrée et génère des métriques de comparaison.
    \item \textbf{ModelSaver} est utilisé par les modèles pour la persistance.
    \item \textbf{PredictionService} utilise les modèles chargés pour effectuer les prédictions.
\end{itemize}

\subsection{Diagramme de séquence}
\label{subsec:sprint5-sequence}

Les diagrammes de séquence suivants illustrent les principaux flux d'exécution du module d'IA.

\subsubsection{Entraînement des modèles}

\begin{figure}[H]
\centering
\includegraphics[scale=0.5]{images/diagramme_sequence_training.png}
\caption{Diagramme de séquence - Entraînement des modèles}
\label{fig:sequence-training}
\end{figure}

\textbf{Description :} Le diagramme de séquence de l'entraînement des modèles montre les interactions entre l'Administrateur, le DataPreprocessor, les cinq modèles, le ModelEvaluator et le ModelSaver. L'administrateur déclenche l'entraînement, le DataPreprocessor prépare les données, chaque modèle est entraîné et évalué, le ModelEvaluator génère le rapport de comparaison, et le ModelSaver sauvegarde les modèles.

Les étapes de l'entraînement des modèles sont :

\begin{enumerate}
    \item L'administrateur déclenche l'entraînement via une interface ou une tâche planifiée.
    \item Le DataPreprocessor se connecte au Data Warehouse pour récupérer les données historiques.
    \item Le DataPreprocessor nettoie les données : suppression des valeurs manquantes, correction des incohérences.
    \item Le DataPreprocessor normalise les données : mise à l'échelle des features numériques.
    \item Le DataPreprocessor encode les variables catégorielles (ex: type d'agence).
    \item Le DataPreprocessor divise les données en ensemble d'entraînement (80\%) et de test (20\%).
    \item Pour chaque modèle (Régression Linéaire, Arbre de Décision, Forêt Aléatoire, Gradient Boosting, Réseau de Neurones) :
        \begin{enumerate}
            \item Le modèle est instancié avec ses paramètres par défaut.
            \item Le modèle est entraîné sur l'ensemble d'entraînement.
            \item Le modèle est évalué sur l'ensemble de test.
            \item Les métriques (RMSE, MAE, R²) sont calculées et stockées.
        \end{enumerate}
    \item Le ModelEvaluator génère un rapport de comparaison des modèles.
    \item Le ModelSaver sauvegarde chaque modèle entraîné dans un fichier (pickle/joblib).
    \item Les métriques et les modèles sauvegardés sont stockés pour une utilisation future.
\end{enumerate}

\subsubsection{Évaluation et comparaison des modèles}

\begin{figure}[H]
\centering
\includegraphics[scale=0.5]{images/diagramme_sequence_evaluation.png}
\caption{Diagramme de séquence - Évaluation et comparaison des modèles}
\label{fig:sequence-evaluation}
\end{figure}

\textbf{Description :} Le diagramme de séquence de l'évaluation et comparaison des modèles montre le ModelEvaluator recevant la liste des modèles entraînés, exécutant les prédictions sur l'ensemble de test, calculant les métriques (RMSE, MAE, R²), générant un tableau comparatif et des visualisations, identifiant le meilleur modèle, et sauvegardant le rapport.

Les étapes de l'évaluation et comparaison des modèles sont :

\begin{enumerate}
    \item Le ModelEvaluator reçoit la liste des modèles entraînés et les données de test.
    \item Pour chaque modèle, le ModelEvaluator exécute la prédiction sur l'ensemble de test.
    \item Le ModelEvaluator calcule les métriques : RMSE (Root Mean Square Error), MAE (Mean Absolute Error), R² (Coefficient de Détermination).
    \item Le ModelEvaluator génère un tableau comparatif avec toutes les métriques.
    \item Le ModelEvaluator crée des visualisations : graphique à barres des métriques, graphique de dispersion des prédictions vs valeurs réelles.
    \item Le ModelEvaluator identifie le meilleur modèle selon chaque métrique.
    \item Le rapport de comparaison est sauvegardé et rendu disponible pour l'API.
\end{enumerate}

\subsubsection{Prédiction via API}

\begin{figure}[H]
\centering
\includegraphics[scale=0.5]{images/diagramme_sequence_prediction_api.png}
\caption{Diagramme de séquence - Prédiction via API}
\label{fig:sequence-prediction-api}
\end{figure}

\textbf{Description :} Le diagramme de séquence de la prédiction via API montre les interactions entre l'Utilisateur, le Frontend React, l'API FastAPI, le PredictionService, le Data Warehouse et les modèles sauvegardés. L'utilisateur sélectionne une agence et un modèle, le frontend envoie une requête POST à l'API, le PredictionService charge le modèle, prépare les données, exécute la prédiction et retourne le résultat avec l'intervalle de confiance.

Les étapes de la prédiction via API sont :

\begin{enumerate}
    \item L'utilisateur sélectionne une agence, un modèle et une période dans l'interface React.
    \item Le frontend envoie une requête POST à l'API FastAPI endpoint /api/prediction.
    \item Le PredictionService reçoit la requête et vérifie les paramètres.
    \item Le PredictionService charge le modèle demandé depuis le disque (s'il n'est pas déjà en mémoire).
    \item Le PredictionService récupère les données historiques de l'agence depuis le Data Warehouse.
    \item Le PredictionService prépare les features pour la prédiction (mêmes transformations que lors de l'entraînement).
    \item Le modèle exécute la prédiction sur les données préparées.
    \item Le PredictionService calcule l'intervalle de confiance de la prédiction.
    \item Le PredictionService retourne le résultat avec les métriques associées.
    \item L'API renvoie la réponse JSON au frontend qui affiche le résultat.
\end{enumerate}

\section{Définition du problème de prédiction}
\label{sec:sprint5-definition}

La définition précise du problème de prédiction est essentielle pour orienter le choix des modèles et l'évaluation des résultats.

\subsection{Prédiction des performances futures des agences}

L'objectif principal est de prédire les performances futures des agences pour aider les responsables à anticiper et à prendre des décisions proactives. Les questions métier auxquelles nous cherchons à répondre sont :

\begin{itemize}
    \item Quel sera le montant des dépôts de l'agence X dans les 3 prochains mois ?
    \item Quelle sera l'évolution du PNB de l'agence Y sur le prochain trimestre ?
    \item Quelles agences risquent de voir leurs performances se dégrader dans les mois à venir ?
\end{itemize}

\subsection{Prédiction des dépôts / crédits / revenus}

Les variables cibles identifiées pour la prédiction sont :

\begin{table}[H]
\centering
\begin{tabular}{|p{3cm}|p{4.5cm}|p{6cm}|}
\hline
\textbf{Variable cible} & \textbf{Unité} & \textbf{Description} \\
\hline
Dépôts & TND (millions) & Montant total des dépôts collectés par l'agence sur une période donnée. \\
\hline
Crédits & TND (millions) & Montant total des crédits accordés par l'agence sur une période donnée. \\
\hline
PNB & TND (millions) & Produit Net Bancaire de l'agence sur une période donnée. \\
\hline
Croissance des dépôts & Pourcentage & Taux de variation des dépôts par rapport à la période précédente. \\
\hline
Nombre de comptes & Unité & Nombre total de comptes actifs dans l'agence. \\
\hline
\end{tabular}
\caption{Variables cibles pour la prédiction}
\label{tab:target-variables}
\end{table}

\subsection{Horizons de prédiction}

Plusieurs horizons de prédiction sont considérés :

\begin{itemize}
    \item \textbf{Court terme} : Mois suivant (M+1) - pour un pilotage opérationnel.
    \item \textbf{Moyen terme} : Trimestre suivant (T+1) - pour une planification tactique.
    \item \textbf{Long terme} : Six mois ou un an - pour une planification stratégique.
\end{itemize}

\subsection{Types de modèles}

Les modèles sont classés en deux catégories :

\begin{itemize}
    \item \textbf{Modèles de régression} : Pour prédire des valeurs numériques continues (montant des dépôts, PNB).
    \item \textbf{Modèles de classification} : Pour prédire des catégories (niveau de performance : bon, moyen, faible).
\end{itemize}

\section{Préparation des données pour l'IA}
\label{sec:sprint5-data-preparation}

La préparation des données est une étape cruciale qui détermine la qualité des prédictions.

\subsection{Feature Engineering}

Le feature engineering consiste à créer de nouvelles variables (features) à partir des données existantes pour améliorer la performance des modèles :

\begin{table}[H]
\centering
\begin{tabular}{|p{3.5cm}|p{3.5cm}|p{7cm}|}
\hline
\textbf{Feature} & \textbf{Type} & \textbf{Description} \\
\hline
Montant des dépôts (historique) & Numérique & Montant des dépôts sur les 12 derniers mois. \\
\hline
Taux de croissance des dépôts & Numérique & Variation en pourcentage par rapport à la période précédente. \\
\hline
Nombre de comptes actifs & Numérique & Nombre total de comptes ouverts et actifs. \\
\hline
Taux de pénétration & Numérique & Ratio clients / population de la zone. \\
\hline
Ancienneté de l'agence & Numérique & Nombre d'années depuis l'ouverture de l'agence. \\
\hline
Nombre d'employés & Numérique & Effectif total de l'agence. \\
\hline
Productivité par employé & Numérique & PNB / nombre d'employés. \\
\hline
Saisonnalité & Catégorielle & Mois, trimestre, période de l'année. \\
\hline
Région & Catégorielle & Région géographique de l'agence. \\
\hline
Type d'agence & Catégorielle & Principale, secondaire, spécialisée. \\
\hline
Indicateurs macroéconomiques & Numérique & PIB local, taux d'inflation (si disponibles). \\
\hline
\end{tabular}
\caption{Features créées pour la prédiction}
\label{tab:features}
\end{table}

\begin{figure}[H]
\centering
\includegraphics[scale=0.5]{images/feature_engineering.png}
\caption{Processus de Feature Engineering}
\label{fig:feature-engineering}
\end{figure}

\textbf{Description :} Le processus de feature engineering consiste à créer de nouvelles variables à partir des données existantes : variables numériques (montant des dépôts, taux de croissance, nombre de comptes, ancienneté, nombre d'employés, productivité), variables catégorielles (saisonnalité, région, type d'agence) et indicateurs macroéconomiques.

\subsection{Nettoyage des données}

Le nettoyage des données comprend les opérations suivantes :

\begin{itemize}
    \item \textbf{Traitement des valeurs manquantes} :
        \begin{itemize}
            \item Pour les variables numériques : imputation par la moyenne, la médiane ou la valeur précédente.
            \item Pour les variables catégorielles : imputation par le mode ou création d'une catégorie "Inconnu".
            \item Suppression des lignes avec trop de valeurs manquantes (> 30\%).
        \end{itemize}
    \item \textbf{Détection et traitement des outliers} :
        \begin{itemize}
            \item Identification des valeurs aberrantes à l'aide de la méthode IQR (Intervalle Interquartile).
            \item Winsorisation : remplacement des valeurs extrêmes par le seuil acceptable.
        \end{itemize}
    \item \textbf{Correction des incohérences} :
        \begin{itemize}
            \item Vérification de la cohérence des données (ex: montant des crédits > 0).
            \item Correction des formats (dates, montants en TND).
        \end{itemize}
\end{itemize}

\subsection{Normalisation des données}

La normalisation est essentielle pour les modèles sensibles à l'échelle des variables :

\begin{itemize}
    \item \textbf{Standardisation (Z-score)} : Transformation pour obtenir une moyenne de 0 et un écart-type de 1. Utilisée pour la Régression Linéaire et le Réseau de Neurones.
    \begin{equation}
        x_{standardized} = \frac{x - \mu}{\sigma}
    \end{equation}
    \item \textbf{Normalisation Min-Max} : Mise à l'échelle dans l'intervalle [0, 1]. Utilisée pour les modèles basés sur les distances (ex: KNN, si utilisé).
    \begin{equation}
        x_{normalized} = \frac{x - x_{min}}{x_{max} - x_{min}}
    \end{equation}
    \item \textbf{Encodage des variables catégorielles} :
        \begin{itemize}
            \item One-Hot Encoding : pour les variables sans ordre (ex: type d'agence).
            \item Label Encoding : pour les variables ordinales (ex: niveau de performance).
        \end{itemize}
\end{itemize}

\subsection{Division Train / Test}

La division des données est effectuée de manière à garantir une évaluation fiable :

\begin{itemize}
    \item \textbf{Ratio Train/Test} : 80\% pour l'entraînement, 20\% pour le test.
    \item \textbf{Division temporelle} : Les données sont divisées par ordre chronologique pour respecter la séquence temporelle (pas de données futures dans l'entraînement).
    \item \textbf{Validation croisée} : Utilisation de la validation croisée k-fold (k=5) pour une évaluation robuste.
\end{itemize}

\begin{table}[H]
\centering
\begin{tabular}{|p{5cm}|p{4cm}|p{4cm}|}
\hline
\textbf{Ensemble} & \textbf{Période} & \textbf{Nombre d'observations} \\
\hline
Entraînement (Train) & Janvier 2022 - Décembre 2024 & 3 200 \\
\hline
Validation (Validation) & Janvier 2025 - Mars 2025 & 400 \\
\hline
Test (Test) & Avril 2025 - Juin 2025 & 400 \\
\hline
\end{tabular}
\caption{Division Train / Validation / Test}
\label{tab:train-test-split}
\end{table}

\section{Implémentation de plusieurs modèles avec Python / Scikit-learn}
\label{sec:sprint5-models}

Cinq modèles sont implémentés et comparés pour identifier le plus performant.

\subsection{Modèle 1 : Régression Linéaire (baseline)}
\label{subsec:linear-regression}

La régression linéaire est utilisée comme modèle de référence (baseline).

\begin{itemize}
    \item \textbf{Principe} : Modélise la relation linéaire entre les features et la variable cible.
    \item \textbf{Équation} : \( y = \beta_0 + \beta_1 x_1 + \beta_2 x_2 + ... + \beta_n x_n + \epsilon \)
    \item \textbf{Avantages} : Simple, interprétable, rapide.
    \item \textbf{Inconvénients} : Ne capture pas les relations non-linéaires.
    \item \textbf{Paramètres} : Aucun paramètre majeur à régler.
\end{itemize}

\subsection{Modèle 2 : Arbre de Décision}
\label{subsec:decision-tree}

L'arbre de décision est un modèle non-linéaire basé sur des règles de décision.

\begin{itemize}
    \item \textbf{Principe} : Divise récursivement l'espace des features en régions homogènes.
    \item \textbf{Avantages} : Interprétable, capture les non-linéarités, pas de normalisation nécessaire.
    \item \textbf{Inconvénients} : Risque de surapprentissage, instabilité.
    \item \textbf{Paramètres clés} : max\_depth (profondeur maximale), min\_samples\_split (nombre minimum d'échantillons pour diviser un nœud).
\end{itemize}

\subsection{Modèle 3 : Forêt Aléatoire}
\label{subsec:random-forest}

La forêt aléatoire est une méthode d'ensemble qui combine plusieurs arbres de décision.

\begin{itemize}
    \item \textbf{Principe} : Agrège les prédictions de nombreux arbres entraînés sur des sous-ensembles aléatoires des données.
    \item \textbf{Avantages} : Robuste, réduit le surapprentissage, gère bien les non-linéarités.
    \item \textbf{Inconvénients} : Moins interprétable qu'un seul arbre, plus lent à l'entraînement.
    \item \textbf{Paramètres clés} : n\_estimators (nombre d'arbres), max\_depth, min\_samples\_split.
\end{itemize}

\subsection{Modèle 4 : Gradient Boosting (XGBoost / LightGBM)}
\label{subsec:gradient-boosting}

Le gradient boosting est une méthode d'ensemble séquentielle qui corrige les erreurs des modèles précédents.

\begin{itemize}
    \item \textbf{Principe} : Construit séquentiellement des arbres, chaque nouvel arbre apprenant des erreurs du précédent.
    \item \textbf{Avantages} : Très performant, gère les non-linéarités, capture les interactions complexes.
    \item \textbf{Inconvénients} : Sensible aux hyperparamètres, plus lent à l'entraînement.
    \item \textbf{Paramètres clés} : learning\_rate (taux d'apprentissage), n\_estimators, max\_depth, subsample.
\end{itemize}

\subsection{Modèle 5 : Réseau de Neurones (MLPRegressor)}
\label{subsec:neural-network}

Le réseau de neurones est un modèle d'apprentissage profond capable de capturer des relations très complexes.

\begin{itemize}
    \item \textbf{Principe} : Réseau de neurones feed-forward avec plusieurs couches cachées.
    \item \textbf{Avantages} : Très flexible, peut approximer n'importe quelle fonction.
    \item \textbf{Inconvénients} : Nécessite beaucoup de données, temps d'entraînement long, difficulté de réglage des hyperparamètres.
    \item \textbf{Paramètres clés} : hidden\_layer\_sizes (structure des couches), activation (fonction d'activation), alpha (régularisation), learning\_rate.
\end{itemize}

\subsection{Modèle optionnel : LSTM (pour données temporelles)}
\label{subsec:lstm}

Le LSTM (Long Short-Term Memory) est un type de réseau de neurones récurrent particulièrement adapté aux séries temporelles.

\begin{itemize}
    \item \textbf{Principe} : Capture les dépendances temporelles à long terme dans les données séquentielles.
    \item \textbf{Avantages} : Idéal pour les séries temporelles, capture les motifs saisonniers.
    \item \textbf{Inconvénients} : Nécessite beaucoup de données, complexe à entraîner.
    \item \textbf{Utilisation} : Optionnel, si les données temporelles sont suffisamment abondantes.
\end{itemize}

\begin{figure}[H]
\centering
\includegraphics[scale=0.5]{images/modeles_ia_comparaison.png}
\caption{Comparaison des modèles d'IA implémentés}
\label{fig:modeles-ia}
\end{figure}

\textbf{Description :} La figure présente les cinq modèles d'IA implémentés : Régression Linéaire (baseline), Arbre de Décision, Forêt Aléatoire, Gradient Boosting (XGBoost/LightGBM) et Réseau de Neurones (MLPRegressor). Chaque modèle est décrit avec ses avantages, inconvénients et paramètres clés.

\section{Évaluation et comparaison des modèles}
\label{sec:sprint5-evaluation}

L'évaluation et la comparaison des modèles sont effectuées à l'aide de métriques standards et de visualisations.

\subsection{Métriques d'évaluation}

Trois métriques principales sont utilisées :

\begin{itemize}
    \item \textbf{RMSE (Root Mean Square Error)} : Erreur quadratique moyenne. Sensible aux grandes erreurs.
    \begin{equation}
        RMSE = \sqrt{\frac{1}{n} \sum_{i=1}^{n} (y_i - \hat{y}_i)^2}
    \end{equation}
    
    \item \textbf{MAE (Mean Absolute Error)} : Erreur absolue moyenne. Interprétable dans l'unité de la variable cible.
    \begin{equation}
        MAE = \frac{1}{n} \sum_{i=1}^{n} |y_i - \hat{y}_i|
    \end{equation}
    
    \item \textbf{R² (Coefficient de Détermination)} : Proportion de la variance expliquée par le modèle. Plus proche de 1, meilleur est le modèle.
    \begin{equation}
        R^2 = 1 - \frac{\sum_{i=1}^{n} (y_i - \hat{y}_i)^2}{\sum_{i=1}^{n} (y_i - \bar{y})^2}
    \end{equation}
\end{itemize}

\subsection{Matrice de comparaison}

La comparaison des modèles est présentée sous forme de tableau :

\begin{table}[H]
\centering
\begin{tabular}{|p{4cm}|p{3cm}|p{3cm}|p{3cm}|}
\hline
\textbf{Modèle} & \textbf{RMSE} & \textbf{MAE} & \textbf{R²} \\
\hline
Régression Linéaire & 1.85 M TND & 1.42 M TND & 0.72 \\
\hline
Arbre de Décision & 1.52 M TND & 1.15 M TND & 0.81 \\
\hline
Forêt Aléatoire & \textbf{1.21 M TND} & \textbf{0.92 M TND} & \textbf{0.89} \\
\hline
Gradient Boosting & 1.28 M TND & 0.98 M TND & 0.87 \\
\hline
Réseau de Neurones & 1.45 M TND & 1.10 M TND & 0.83 \\
\hline
\end{tabular}
\caption{Comparaison des performances des modèles}
\label{tab:model-comparison}
\end{table}

\subsection{Visualisation des comparaisons}

Les comparaisons sont visualisées à l'aide de :

\begin{itemize}
    \item \textbf{Graphique à barres} : Comparaison des métriques (RMSE, MAE, R²) pour chaque modèle.
    \item \textbf{Graphique de dispersion} : Valeurs prédites vs valeurs réelles pour visualiser la qualité des prédictions.
    \item \textbf{Boxplot des erreurs} : Distribution des erreurs de prédiction pour chaque modèle.
    \item \textbf{Heatmap des corrélations} : Corrélation entre les features et la variable cible.
\end{itemize}

\begin{figure}[H]
\centering
\includegraphics[scale=0.5]{images/model_comparison_chart.png}
\caption{Graphique de comparaison des modèles}
\label{fig:model-comparison-chart}
\end{figure}

\textbf{Description :} Le graphique de comparaison des modèles présente les métriques (RMSE, MAE, R²) pour chacun des cinq modèles. La Forêt Aléatoire se distingue avec le meilleur R² (0.89) et le plus faible RMSE (1.21 M TND), tandis que la Régression Linéaire constitue la baseline avec des performances plus faibles.

\subsection{Sélection du meilleur modèle}

Le modèle de Forêt Aléatoire est sélectionné comme modèle principal pour les raisons suivantes :

\begin{itemize}
    \item \textbf{Meilleur R²} : 0.89, expliquant 89\% de la variance des données.
    \item \textbf{Plus faible RMSE} : 1.21 M TND, indiquant des prédictions plus précises.
    \item \textbf{Robustesse} : Moins sensible aux outliers que d'autres modèles.
    \item \textbf{Compromis performance/complexité} : Bonne performance sans être excessivement complexe à entraîner.
\end{itemize}

\section{Sauvegarde des modèles entraînés (pickle / joblib)}
\label{sec:sprint5-save}

La sauvegarde des modèles entraînés permet de les réutiliser sans avoir à les ré-entraîner.

\subsection{Format de sauvegarde}

Deux formats sont utilisés :

\begin{itemize}
    \item \textbf{Pickle} : Format standard de sérialisation Python. Permet de sauvegarder n'importe quel objet Python.
    \item \textbf{Joblib} : Format optimisé pour les grands tableaux NumPy, plus efficace pour les modèles Scikit-learn.
\end{itemize}

\subsection{Structure de sauvegarde}

La sauvegarde inclut :

\begin{itemize}
    \item \textbf{Le modèle entraîné} : L'objet modèle Scikit-learn (ou autre) sérialisé.
    \item \textbf{Les métadonnées} : Informations sur le modèle (nom, date d'entraînement, métriques de performance, version).
    \item \textbf{Les paramètres de normalisation} : Moyennes, écarts-types, minimums, maximums utilisés pour la normalisation.
    \item \textbf{Les features utilisées} : Liste des features (colonnes) utilisées pour l'entraînement.
    \item \textbf{La version du modèle} : Numéro de version pour suivre les évolutions.
\end{itemize}

\subsection{Organisation des fichiers}

Les modèles sont organisés dans une structure de dossiers :

\begin{verbatim}
models/
├── linear_regression/
│   ├── model.pkl
│   ├── metadata.json
│   └── scaler.pkl
├── decision_tree/
│   ├── model.pkl
│   ├── metadata.json
│   └── scaler.pkl
├── random_forest/
│   ├── model.pkl
│   ├── metadata.json
│   └── scaler.pkl
├── gradient_boosting/
│   ├── model.pkl
│   ├── metadata.json
│   └── scaler.pkl
└── neural_network/
    ├── model.pkl
    ├── metadata.json
    └── scaler.pkl
\end{verbatim}

\subsection{Gestion des versions}

Un système de versioning permet de garder trace des évolutions :

\begin{itemize}
    \item \textbf{Version sémantique} : MAJEUR.MINEUR.PATCH (ex: 1.0.0, 1.1.0, 2.0.0).
    \item \textbf{Date d'entraînement} : Horodatage de l'entraînement.
    \item \textbf{Performance} : Métriques RMSE, MAE, R² pour chaque version.
    \item \textbf{Archivage} : Possibilité de conserver plusieurs versions pour comparaison.
\end{itemize}

\begin{figure}[H]
\centering
\includegraphics[scale=0.5]{images/sauvegarde_modeles.png}
\caption{Structure de sauvegarde des modèles}
\label{fig:sauvegarde-modeles}
\end{figure}

\textbf{Description :} La structure de sauvegarde des modèles organise chaque modèle dans un dossier dédié contenant le modèle sérialisé (model.pkl), les métadonnées (metadata.json) et le scaler de normalisation (scaler.pkl). Un système de versioning sémantique (MAJEUR.MINEUR.PATCH) permet de suivre les évolutions.

\section{Création d'un endpoint FastAPI pour la prédiction}
\label{sec:sprint5-endpoint}

L'endpoint FastAPI permet d'exposer les modèles de prédiction via l'API.

\subsection{Endpoint /predict}

\begin{itemize}
    \item \textbf{URL} : POST /api/prediction
    \item \textbf{Authentification} : Requise (JWT)
    \item \textbf{Corps de la requête} :
        \begin{itemize}
            \item agence\_id (int) : Identifiant de l'agence.
            \item modele (string) : Nom du modèle à utiliser (linear\_regression, decision\_tree, random\_forest, gradient\_boosting, neural\_network).
            \item periode (string) : Période de prédiction (next\_month, next\_quarter).
            \item indicateur (string) : Indicateur à prédire (depots, credits, pnb).
        \end{itemize}
    \item \textbf{Réponse} :
        \begin{itemize}
            \item prediction (float) : Valeur prédite.
            \item confidence\_interval (array) : Intervalle de confiance [borne\_inf, borne\_sup].
            \item metrics (object) : Métriques du modèle (RMSE, MAE, R²).
            \item model\_name (string) : Nom du modèle utilisé.
            \item timestamp (string) : Date et heure de la prédiction.
        \end{itemize}
\end{itemize}

\subsection{Exemple de réponse}

\begin{verbatim}
{
    "prediction": 12.45,
    "confidence_interval": [11.85, 13.05],
    "metrics": {
        "rmse": 1.21,
        "mae": 0.92,
        "r2": 0.89
    },
    "model_name": "random_forest",
    "target": "depots",
    "agence_id": 1,
    "periode": "next_month",
    "timestamp": "2025-01-15T14:30:00Z"
}
\end{verbatim}

\subsection{Endpoint /modeles}

Un endpoint supplémentaire permet de lister les modèles disponibles :

\begin{itemize}
    \item \textbf{URL} : GET /api/prediction/modeles
    \item \textbf{Authentification} : Requise (JWT)
    \item \textbf{Réponse} : Liste des modèles avec leurs métriques et leur disponibilité.
\end{itemize}

\subsection{Gestion des erreurs}

L'endpoint gère les cas d'erreur suivants :

\begin{itemize}
    \item \textbf{Modèle non trouvé} : 404 Not Found.
    \item \textbf{Agence non trouvée} : 404 Not Found.
    \item \textbf{Données insuffisantes} : 400 Bad Request.
    \item \textbf{Erreur de prédiction} : 500 Internal Server Error.
    \item \textbf{Non autorisé} : 403 Forbidden (si l'utilisateur n'a pas accès à l'agence).
\end{itemize}

\begin{figure}[H]
\centering
\includegraphics[scale=0.5]{images/endpoint_prediction_api.png}
\caption{Schéma de l'endpoint /api/prediction}
\label{fig:endpoint-prediction}
\end{figure}

\textbf{Description :} L'endpoint /api/prediction reçoit une requête POST avec les paramètres (agence\_id, modele, periode, indicateur), vérifie l'authentification JWT, appelle le PredictionService qui charge le modèle et exécute la prédiction, puis retourne une réponse JSON avec la valeur prédite, l'intervalle de confiance et les métriques du modèle.

\section{Comparaison des modèles côté React}
\label{sec:sprint5-react-comparison}

L'interface React permet de visualiser et comparer les performances des modèles.

\subsection{Sélection du modèle via dropdown}

L'interface offre un menu déroulant (dropdown) pour sélectionner :

\begin{itemize}
    \item \textbf{Modèle} : Choisir entre Régression Linéaire, Arbre de Décision, Forêt Aléatoire, Gradient Boosting, Réseau de Neurones.
    \item \textbf{Agence} : Sélectionner l'agence pour laquelle la prédiction est demandée.
    \item \textbf{Période} : Choisir l'horizon de prédiction (mois suivant, trimestre suivant).
    \item \textbf{Indicateur} : Choisir la variable à prédire (dépôts, crédits, PNB).
\end{itemize}

\subsection{Affichage des résultats}

Les résultats sont affichés sous différentes formes :

\begin{itemize}
    \item \textbf{Carte de prédiction} : Affichage de la valeur prédite avec l'intervalle de confiance.
    \item \textbf{Métriques du modèle} : Affichage des métriques (RMSE, MAE, R²) pour évaluer la fiabilité.
    \item \textbf{Graphique comparatif} : Visualisation des prédictions de plusieurs modèles côte à côte.
\end{itemize}

\subsection{Comparaison des métriques}

Un tableau comparatif permet de visualiser les métriques de tous les modèles :

\begin{table}[H]
\centering
\begin{tabular}{|p{3.5cm}|p{2.5cm}|p{2.5cm}|p{2.5cm}|}
\hline
\textbf{Modèle} & \textbf{RMSE} & \textbf{MAE} & \textbf{R²} \\
\hline
Régression Linéaire & 1.85 M TND & 1.42 M TND & 0.72 \\
\hline
Arbre de Décision & 1.52 M TND & 1.15 M TND & 0.81 \\
\hline
Forêt Aléatoire & \textbf{1.21 M TND} & \textbf{0.92 M TND} & \textbf{0.89} \\
\hline
Gradient Boosting & 1.28 M TND & 0.98 M TND & 0.87 \\
\hline
Réseau de Neurones & 1.45 M TND & 1.10 M TND & 0.83 \\
\hline
\end{tabular}
\caption{Comparaison des métriques - Interface React}
\label{tab:react-model-comparison}
\end{table}

\section{Interface React pour la prédiction}
\label{sec:sprint5-react-interface}

L'interface React dédiée à la prédiction offre une expérience utilisateur intuitive.

\subsection{Sélection de l'agence}

\begin{itemize}
    \item \textbf{Sélecteur d'agence} : Menu déroulant avec la liste de toutes les agences (filtrée par rôle).
    \item \textbf{Recherche rapide} : Champ de recherche pour trouver rapidement une agence.
    \item \textbf{Agence par défaut} : Pour un responsable d'agence, son agence est présélectionnée.
\end{itemize}

\subsection{Sélection du modèle}

\begin{itemize}
    \item \textbf{Boutons de sélection} : Choix du modèle via des boutons ou un menu déroulant.
    \item \textbf{Option "Tous les modèles"} : Visualisation simultanée de toutes les prédictions.
    \item \textbf{Indicateur du meilleur modèle} : Surbrillance du modèle avec le meilleur R².
\end{itemize}

\subsection{Visualisation des prédictions}

Les prédictions sont visualisées sous plusieurs formes :

\begin{itemize}
    \item \textbf{Jauge (Gauge)} : Visualisation de la valeur prédite par rapport aux valeurs historiques et aux objectifs.
    \item \textbf{Graphique d'évolution} : Projection de la tendance future sur le graphique historique.
    \item \textbf{Carte de prédiction} : Affichage de la valeur prédite avec l'intervalle de confiance.
    \item \textbf{Score de confiance} : Indicateur de la fiabilité de la prédiction basé sur l'intervalle de confiance et les métriques du modèle.
    \item \textbf{Tableau comparatif} : Visualisation côte à côte des prédictions de tous les modèles.
\end{itemize}

\begin{figure}[H]
\centering
\includegraphics[scale=0.5]{images/maquette_interface_prediction_react.png}
\caption{Maquette de l'interface de prédiction React}
\label{fig:react-prediction-interface}
\end{figure}

\textbf{Description :} La maquette de l'interface de prédiction React présente un sélecteur d'agence, des boutons de sélection des modèles (Régression Linéaire, Arbre de Décision, Forêt Aléatoire, Gradient Boosting, Réseau de Neurones), une jauge de visualisation, une carte de prédiction avec intervalle de confiance, et un tableau comparatif des métriques (RMSE, MAE, R²).

\subsection{Fonctionnalités avancées}

\begin{itemize}
    \item \textbf{Export des résultats} : Téléchargement des prédictions en PDF ou Excel.
    \item \textbf{Historique des prédictions} : Visualisation des prédictions précédentes pour suivi de la performance du modèle.
    \item \textbf{Comparaison avec la réalité} : Affichage des valeurs réelles une fois disponibles pour évaluer la précision du modèle.
    \item \textbf{Scénarios "what-if"} : Simulation de l'impact de modifications des paramètres sur les prédictions.
\end{itemize}

\section{Tests et validation}
\label{sec:sprint5-tests}

Les tests et la validation sont essentiels pour garantir la fiabilité du module de prédiction.

\subsection{Tests unitaires}

Les tests unitaires couvrent :

\begin{itemize}
    \item \textbf{DataPreprocessor} : Validation des fonctions de nettoyage, normalisation et encodage.
    \item \textbf{Chaque modèle} : Vérification de l'entraînement, de la prédiction et de l'évaluation.
    \item \textbf{ModelEvaluator} : Vérification du calcul des métriques.
    \item \textbf{ModelSaver} : Vérification de la sauvegarde et du chargement.
\end{itemize}

\subsection{Tests de validation croisée}

La validation croisée k-fold (k=5) est utilisée pour :

\begin{itemize}
    \item Évaluer la stabilité des modèles sur différents sous-ensembles de données.
    \item Détecter le surapprentissage.
    \item Sélectionner les hyperparamètres optimaux.
\end{itemize}

\subsection{Tests sur données réelles}

Les modèles sont testés sur des données réelles historiques :

\begin{itemize}
    \item \textbf{Backtesting} : Les modèles sont testés sur des données historiques qu'ils n'ont pas vues.
    \item \textbf{Comparaison avec la réalité} : Les prédictions sont comparées aux valeurs réelles pour évaluer la précision.
    \item \textbf{Suivi de performance} : Les métriques sont suivies dans le temps pour détecter une éventuelle dégradation.
\end{itemize}

\subsection{Résultats de validation}

\begin{table}[H]
\centering
\begin{tabular}{|p{4cm}|p{3cm}|p{3cm}|p{3cm}|}
\hline
\textbf{Modèle} & \textbf{RMSE (test)} & \textbf{R² (test)} & \textbf{Validation croisée} \\
\hline
Régression Linéaire & 1.85 M TND & 0.72 & 0.71 ± 0.02 \\
\hline
Arbre de Décision & 1.52 M TND & 0.81 & 0.79 ± 0.04 \\
\hline
Forêt Aléatoire & \textbf{1.21 M TND} & \textbf{0.89} & \textbf{0.88 ± 0.02} \\
\hline
Gradient Boosting & 1.28 M TND & 0.87 & 0.86 ± 0.03 \\
\hline
Réseau de Neurones & 1.45 M TND & 0.83 & 0.82 ± 0.05 \\
\hline
\end{tabular}
\caption{Résultats de validation des modèles}
\label{tab:validation-results}
\end{table}

\section{Livrables du Sprint}
\label{sec:sprint5-livrables}

À l'issue de ce sprint, les livrables suivants sont produits :

\begin{table}[H]
\centering
\begin{tabular}{|p{4cm}|p{9cm}|}
\hline
\textbf{Livrable} & \textbf{Description} \\
\hline
Modèles entraînés & 5 modèles de prédiction entraînés (Régression Linéaire, Arbre de Décision, Forêt Aléatoire, Gradient Boosting, Réseau de Neurones). \\
\hline
Rapport de comparaison & Document détaillant les performances de chaque modèle avec métriques et graphiques. \\
\hline
Modèles sauvegardés & Fichiers pickle/joblib des modèles entraînés avec leurs métadonnées. \\
\hline
Endpoint FastAPI & Endpoint /api/prediction et /api/prediction/modeles pour les prédictions. \\
\hline
Interface React & Interface de sélection et visualisation des prédictions. \\
\hline
Documentation technique & Guide d'utilisation des modèles et de l'API de prédiction. \\
\hline
Rapport de validation & Résultats des tests de validation croisée et des backtests. \\
\hline
\end{tabular}
\caption{Livrables du sprint 5}
\label{tab:sprint5-livrables}
\end{table}



\textbf{Description :} Les livrables du sprint 5 comprennent : les 5 modèles entraînés, le rapport de comparaison des performances, les modèles sauvegardés (fichiers pickle/joblib), les endpoints FastAPI (/api/prediction et /api/prediction/modeles), l'interface React de sélection et visualisation, la documentation technique et le rapport de validation.

\section*{Conclusion}
\addcontentsline{toc}{section}{Conclusion}

Ce cinquième sprint a permis de développer un module d'intelligence artificielle complet pour la prédiction des performances des agences de la BTL. Cinq modèles différents ont été implémentés, évalués et comparés, avec la Forêt Aléatoire comme modèle le plus performant (R² = 0.89, RMSE = 1.21 M TND).

La préparation rigoureuse des données (feature engineering, nettoyage, normalisation) a été essentielle pour garantir la qualité des prédictions. L'évaluation et la comparaison des modèles ont permis d'identifier le meilleur algorithme pour notre problème spécifique, tout en offrant une transparence totale sur les performances de chaque approche.

La sauvegarde des modèles et l'intégration dans l'API FastAPI permettent désormais d'exposer les prédictions de manière sécurisée et performante. L'interface React développée offre aux responsables de la BTL un outil intuitif pour visualiser les prédictions, comparer les modèles et prendre des décisions éclairées.

Ce module d'IA constitue une valeur ajoutée significative pour la BTL, permettant d'anticiper les performances futures, de détecter les risques potentiels et d'optimiser les stratégies de développement des agences. Le prochain sprint sera dédié au développement du chatbot intelligent, qui utilisera ces modèles de prédiction pour répondre aux questions des utilisateurs en langage naturel.